% ==============================================================================
% 07_appendix.tex – Appendix (included after the bibliography)
% ==============================================================================
\appendix
\chapter{Appendix}
\label{ch:anhang}

\section{Structure of the Code}
\label{sec:anhang_code}
The code is written in Python with PyTorch and torchvision. The most important
parts are:
\begin{itemize}[nosep]
  \item \code{main.py}: loads GazeCapture, normalizes the targets and creates
        the splits;
  \item \code{model.py}, \code{dataset.py}: model, losses, training, data
        loading and stored features;
  \item \code{metrics.py}: diagonal error and baselines (one single source for
        all scripts);
  \item \code{exp.py}: hyperparameter study with five repeats;
  \item \code{train\_base.py}, \code{personal\_exp.py}: base models and
        personalization (E0--E3);
  \item \code{moving\_average.py}: temporal smoothing of the stored
        predictions (E4);
  \item \code{session\_design.py}: coverage, transition frames and cleaned
        labels;
  \item \code{spatial\_analysis.py}, \code{plot\_targets.py}: figures;
  \item \code{auswertung.py}: creates the tables from the result files;
  \item notebook \code{notebooks/Masterarbeit.ipynb}: runs all experiments and
        creates tables and figures;
  \item \code{demo/}: web demo (server and browser interface).
\end{itemize}

\section{Commands to Repeat the Experiments}
\label{sec:anhang_befehle}
All commands are run in the main folder of the project (see also
\code{README.md}).

\paragraph{Prepare the data.}
\begin{verbatim}
python main.py
\end{verbatim}

\paragraph{Baselines and hyperparameter study (GazeCapture).}
\begin{verbatim}
python exp.py --experiment baseline
python exp.py --experiment all --sample-seed 0 --val-size 200
python exp.py --experiment trainable_layers --include-resnet
\end{verbatim}

\paragraph{Spatial analysis and figures.}
\begin{verbatim}
python spatial_analysis.py --input results/heatmap/base_loss1.csv \
       --compare results/heatmap/base_loss2.csv --out results/figures
python plot_targets.py --out results/figures
\end{verbatim}

\paragraph{Personalization (E0--E3).}
\begin{verbatim}
python train_base.py
python personal_exp.py --sessions all --base norm_subject norm_random
python personal_exp.py --sessions gen4 --labels clean
python personal_exp.py --sessions 21 22 --split random --seeds 42 43 44
python personal_exp.py --sessions 21 22 --split sequential --seeds 42 43 44
python personal_exp.py --split two_videos --pairs 21:22,22:21 --seeds 42 43 44
python personal_exp.py --sessions all --split sequential
python personal_exp.py --summary
\end{verbatim}

\paragraph{Temporal smoothing (E4).}
\begin{verbatim}
python moving_average.py
\end{verbatim}

\paragraph{Session design and cleaned labels.}
\begin{verbatim}
python session_design.py
python session_design.py --write-clean
\end{verbatim}

It is easier to use the notebook \code{notebooks/Masterarbeit.ipynb}: cell~1 trains the
base models (E0), cell~3 runs the personalization (E1 to E3) and creates all
tables.
